\subsection{绝对收敛积分}

定义 2. 我们说规则函数 \(\mathbf{f}\) 在区间 \(\mathrm{I} \subset \mathbf{R}\) 上的积分是绝对收敛的，如果正函数 \(\| \mathbf{f}(x) \|\) 的积分收敛。

命题 4. 若 \(\mathbf{f}\) 在 \(\mathrm{I}\) 上的积分绝对收敛，则它收敛，且有

\[
\left\| \int_ {I} \mathbf {f} (t) d t \right\| \leqslant \int_ {I} \| \mathbf {f} (t) \| d t.\tag{2}
\]

事实上，对每个紧区间 \(J \subset I\) 有（II, p.~61, 公式 (16)）

\[
\left\| \int_ {J} \mathbf {f} (t) d t \right\| \leqslant \int_ {J} \| \mathbf {f} (t) \| d t.\tag{3}
\]

若正函数 \(\|\mathbf{f}(x)\|\) 的积分收敛，则对每个 \(\varepsilon > 0\) 存在一个含于 I 的紧区间 \([\alpha, \beta]\)，使得对每个紧区间 \([x, y]\)（含于 I 且与 \([\alpha, \beta]\) 无公共内点），有 \(\int_{x}^{y} \|\mathbf{f}(t) dt\| \leqslant \varepsilon\)（II, p.~65, 命题 1）；由此推出 \(\left\|\int_{x}^{y} \mathbf{f}(t) dt\right\| \leqslant \varepsilon\)，这表明 I 上的积分收敛（II, p.~16, 命题 1）；在 (3) 中取极限便得不等式 (2)。

推论。设 E, F, G 是 R 上的三个完备赋范空间，\((\mathbf{x}, \mathbf{y}) \mapsto [\mathbf{x}, \mathbf{y}]\) 是 \(E \times F\) 到 G 中的连续双线性映射。设 f, g 是 I 上的两个规则函数，分别取值于 E 与 F 中。若 f 在 I 上有界且 \(\mathbf{g}\) 的积分在 I 上绝对收敛，则 \([\mathbf{f},\mathbf{g}]\) 的积分绝对收敛。

事实上，存在一个数 \(h > 0\)，使得 \(\| [\mathbf{x},\mathbf{y}] \| \leqslant h \| \mathbf{x} \|\) . \(\| \mathbf{y} \|\) 恒成立（Gen.~Top., IX, p.~173, 定理 1）；若令 \(k = \sup_{x \in I} \| \mathbf{f}(x) \|\)，则在 I 上有 \(\| [\mathbf{f}(x),\mathbf{g}(x)] \| \leqslant hk \| \mathbf{g}(x) \|\)；于是比较原理表明 {[}f,g{]} 的积分绝对收敛，且由 (2)，

\[
\left\| \int_ {I} [ \mathbf {f} (t). \mathbf {g} (t) ] d t \right\| \leqslant h k \int_ {I} \| \mathbf {g} (t) \| d t.
\]

注记。一个积分可以收敛而非绝对收敛；II, p.~64 的例 3 就表明了这一点，那里通项为 \(u_{n}\) 的级数收敛而非绝对收敛。

